\scanpage{38}

Rappelons

\[
\underset{\substack{\Vert\\p_3^{-1}(\alpha)}}{E(\unsure{\sigma_*|\alpha_*|\alpha})}\subset\Pi_\infty\xrightarrow{p_3}\Pi_3
\tag*{(42)}
\]

\[
\underset{\deleted{\unsure{-\bar j}}}{\unsure{\sigma|\alpha}}
=(\bar j,\deleted{\bar{\vphantom j}}j,1)\bmod\mu_3
=\underset{\substack{\rotatebox{-90}{$\in$}\\\Pi_3}}{\{(\bar j,\deleted{\bar{\vphantom j}}j,1),(\deleted{\bar{\vphantom j}}j,1,\bar j),(1,\bar j,\deleted{\bar{\vphantom j}}j)\}}
\]

\deleted{Posons} Nous \deleted{prendrons} utiliserons

\[
(\qquad j=\exp\tfrac{2i\pi}{3},\qquad \bar j=\exp-\tfrac{2i\pi}{3}
\]

et poserons

\[
\zeta_n=\exp\tfrac{2i\pi}{3n}
\tag*{(43)}
\]

et

\[
\left\{\begin{aligned}
\deleted{\unsure{\alpha_n\text{ soit }\zeta_n}}\ \sigma_n
&=(\bar\zeta_n,\deleted{\bar{\vphantom\zeta}}\zeta_n\deleted{j},1\deleted{\unsure{\zeta}})
&&\left\{\begin{aligned}
&\text{considéré soit comme un}\\[-.3ex]
&\text{élément de }\overline C_n\text{ au dessus de }\sigma=-j\\
&\deleted{\bmod\mu_n\text{ dans}}\\
&\text{soit comme élément de }\mu_{3n}^{\,3}\bmod\mu_{\unsure{3n|3}}\\
&\deleted{\unsure{\text{définisse}}}\\[-.3ex]
&=\Pi_{3n}
\end{aligned}\right.\\
\sigma_\infty&=\begin{aligned}
&\deleted{\unsure{(\zeta_\infty,\bar\zeta_\infty,1)}}\\
&(\bar\zeta_\infty,\zeta_\infty,1)
\end{aligned}
&&\left\{\begin{aligned}
&\text{considéré soit comme pts de}\\
&\overline C_\infty\text{ au dessus de }\alpha=-j\\
&\text{soit comme él. de }\mu_\infty^{\,3}\bmod\mu_\infty\\
&=\Pi_\infty
\end{aligned}\right.
\end{aligned}\right.
\]

où

\[
\zeta_\infty=(\zeta_n)\in T
\tag*{(\deleted{\unsure{4}}44)}
\]

est la base canonique de $T$, donnée par l'exponentielle complexe, permettant donc de définir un \unsure{isom. can|iso. can.}

\[
\widehat{\mathbb Z}\mathrel{\overset{\varphi}{\simeq}}T .
\tag*{(45)}
\]

% archive sheet number: 37
